\newpage
\section{Aufgaben mit Schwerpunkt Systeme gewöhnlicher Differentialgleichungen. (Kapitel 5)}

\section*{Aufgabe 1 (12 Punkte): Kurzfragen}
Beantworten Sie die folgenden Fragen kurz und präzise.

\begin{enumerate}
    \item Was ist der Unterschied zwischen der \textit{expliziten} und der \textit{impliziten} Form einer Differentialgleichung $n$-ter Ordnung?
    \item Definieren Sie ein \textit{System von $n$ Differentialgleichungen 1. Ordnung}.
    \item Was versteht man unter dem \textit{Graphen} einer Lösung $y(t)$ eines DGL-Systems?
    \item Definieren Sie eine \textit{Anfangswertaufgabe} (AWA) für ein System 1. Ordnung.
    \item Wann spricht man von einer \textit{globalen Lösung} einer Anfangswertaufgabe?
    \item Formulieren Sie den Satz über die Äquivalenz einer AWA zu einer \textit{Integralgleichung} (Satz 5.1.4).
    \item Was ist eine \textit{autonome} Differentialgleichung?
    \item Wie transformiert man eine skalare DGL 2. Ordnung $y'' = f(t, y, y')$ in ein System 1. Ordnung?
    \item Welches Kriterium an die Funktion $f(t, y)$ garantiert (nach Picard-Lindelöf) die \textit{eindeutige} lokale Lösbarkeit einer AWA?
    \item Was ist die \textit{Lipschitz-Bedingung} bzgl. der zweiten Variablen?
    \item Erklären Sie kurz das Prinzip der \textit{Picard-Iteration}.
    \item Was bedeutet \textit{lineare Unabhängigkeit} von Lösungen eines homogenen linearen DGL-Systems?
\end{enumerate}

\newpage

\section*{Aufgabe 2: Umwandlung in Systeme 1. Ordnung (8 Punkte)}
Gegeben sei die lineare Differentialgleichung 3. Ordnung:
\[ y''' - 2y'' + y' - 5y = \sin(t) \]
Wandeln Sie diese Gleichung in ein System von Differentialgleichungen 1. Ordnung der Form $z' = A z + b(t)$ um. Geben Sie die Matrix $A$ und den Vektor $b(t)$ explizit an.

\newpage

\section*{Aufgabe 3: AWA und Integralgleichung (10 Punkte)}
Betrachten Sie die Anfangswertaufgabe:
\[ y' = 2t \cdot y, \quad y(0) = 1 \]
\begin{enumerate}
    \item Geben Sie die zugehörige Integralgleichung gemäß Satz 5.1.4 an.
    \item Berechnen Sie die ersten zwei Iterationen $y_1(t)$ und $y_2(t)$ der Picard-Iteration, ausgehend vom Startwert $y_0(t) = 1$.
\end{enumerate}

\newpage

\section*{Aufgabe 4: Lineare Systeme mit konstanten Koeffizienten (10 Punkte)}
Bestimmen Sie die allgemeine Lösung des folgenden homogenen Systems 1. Ordnung:
\[
\begin{pmatrix} y_1' \\ y_2' \end{pmatrix} = \begin{pmatrix} 1 & 2 \\ 2 & 1 \end{pmatrix} \begin{pmatrix} y_1 \\ y_2 \end{pmatrix}
\]
\textit{Hinweis: Berechnen Sie die Eigenwerte und Eigenvektoren der Systemmatrix.}

\newpage

\section*{Aufgabe 5: Trennung der Variablen (10 Punkte)}
Lösen Sie die folgende Anfangswertaufgabe durch Trennung der Variablen:
\[ y' = \frac{e^t}{y}, \quad y(0) = \sqrt{2} \]
Geben Sie den maximalen Definitionsbereich der Lösung an.